\chapter{Implementation}

\section{Software Organisation}

The implementation is organised as a modular MATLAB/Simulink workspace. The
\texttt{data} directory contains reusable aircraft, atmosphere, aerodynamic,
propulsion and actuator data. The \texttt{plant} directory contains adapters
that convert those data into forces, moments and state derivatives. The
\texttt{tests} directory contains feature-oriented MATLAB checks, while
\texttt{tools} stores scripts used to reproduce DATCOM extraction and validation
artifacts.

The main implementation rule is separation of concerns. Fixed aircraft
properties, environment models, aerodynamic data, actuator dynamics and plant
force build-up are kept in separate modules. This keeps each assumption
traceable and allows individual layers to be replaced without changing the full
plant interface.

\section{Implemented Modules}

The current data layer contains:
\begin{itemize}
    \item geometry, mass, inertia, constants and ISA atmosphere modules,
    \item a staged aerodynamic database with CRM clean-cruise data, DATCOM-ready
    derivative layers and high-lift increments,
    \item a B777-like propulsion seed with thrust lapse and spool dynamics,
    \item a control-surface actuator seed with position and rate limits,
    \item reference cases, configurations and modelling conventions.
\end{itemize}

The current plant-facing adapters are:
\begin{itemize}
    \item \texttt{b777\_aero\_forces\_moments.m}, which converts aerodynamic
    coefficients into body-axis forces and moments,
    \item \texttt{b777\_engine\_forces\_moments.m}, which converts left and
    right engine thrust into body-axis propulsion loads,
    \item \texttt{b777\_actuator\_dynamics.m}, which maps commanded
    control-surface deflections to actual surface-state derivatives.
\end{itemize}

\section{Aerodynamic Data Processing}

The aerodynamic database is built from versioned source files and reproducible
helper scripts. NASA CRM data provide the active clean-cruise longitudinal
baseline. Digital DATCOM outputs are parsed into candidate derivative tables
for static, damping and control effects. Derived gap-fill candidates cover
interface terms that are not printed directly by the current DATCOM workflow.
Trim, modal and promotion-gate reports are stored as validation evidence before
candidate data are promoted.

Generated DATCOM run folders are treated as reproducible workspace outputs and
are not stored as source data. The committed repository retains the input
decks, candidate CSV files, reports and validation artifacts needed to audit
the modelling chain.

\section{Feature Tests}

Feature tests check module interfaces and modelling conventions rather than
certifying aircraft fidelity. They currently verify:
\begin{itemize}
    \item project structure and MATLAB path assumptions,
    \item state, input, unit and sign conventions,
    \item geometry, mass, inertia, atmosphere and reference-case consistency,
    \item aerodynamic source inventory, coefficient behaviour and load signs,
    \item engine thrust mapping, thrust lapse, spool derivatives and asymmetric
    thrust moments,
    \item actuator position limits, rate limits and command-to-state mapping.
\end{itemize}

These tests are intended to catch convention errors early. Physical validation
is handled separately through trim, modal and later simulation studies.

\section{Current Status}

At the current stage the project contains the reusable data layer, aerodynamic
load adapter, propulsion load adapter and control-surface actuator adapter. The
next implementation step is to assemble the total force and moment build-up and
integrate the nonlinear six-degree-of-freedom equations in MATLAB/Simulink.
